\section{Bài toán năng lượng của phản ứng}

\subsection{Hiệu suất sử dụng năng lượng}

Trong thực tế, không phải toàn bộ nhiệt lượng do phản ứng tỏa ra đều được chuyển thành phần năng lượng có ích. Một phần bị thất thoát ra môi trường. Nếu hiệu suất sử dụng năng lượng là \(H\%\), ta có:
\[
Q_{\text{ích}}=Q_{\text{tỏa}}\frac{H}{100}.
\]
Do đó:
\[
Q_{\text{tỏa}}=\frac{Q_{\text{ích}}}{H/100}.
\]

\begin{ghinho}{Phân biệt hai loại hiệu suất}
\begin{itemize}
\item \textbf{Hiệu suất phản ứng} cho biết mức độ chuyển hóa chất đầu thành sản phẩm.
\item \textbf{Hiệu suất sử dụng năng lượng} cho biết phần năng lượng tỏa ra thực sự được dùng cho mục đích mong muốn.
\end{itemize}
Hai đại lượng này có thể cùng xuất hiện trong một bài toán nhưng không được dùng thay cho nhau.
\end{ghinho}

\begin{vidu}
Một bếp gas cần cung cấp \SI{18}{\mega\joule} nhiệt có ích để đun nấu. Hiệu suất sử dụng nhiệt của bếp là \(60\%\). Tính nhiệt lượng mà nhiên liệu phải tỏa ra.
\end{vidu}
\begin{loigiai}
\[
Q_{\text{tỏa}}=\frac{18}{0.60}=\SI{30}{\mega\joule}.
\]

\textbf{Vậy,} nhiên liệu phải tỏa ra \SI{30}{\mega\joule}.
\end{loigiai}

\subsection{Từ năng lượng cần thiết đến lượng nhiên liệu}

Các bài tuyển sinh thường yêu cầu đi theo chuỗi:
\[
\text{năng lượng có ích}
\longrightarrow
\text{năng lượng nhiên liệu phải tỏa}
\longrightarrow
\text{số mol hoặc khối lượng nhiên liệu}.
\]

\begin{vidu}
Một lò cần \SI{504}{\mega\joule} nhiệt có ích. Lò đốt methane, biết \(1\) mol \ce{CH4} cháy hoàn toàn tỏa \SI{890}{\kilo\joule}; hiệu suất sử dụng nhiệt của lò là \(70\%\). Tính khối lượng methane cần dùng. Cho \(M_{\ce{CH4}}=16\) g/mol.
\end{vidu}
\begin{loigiai}
Nhiệt lượng methane phải tỏa:
\[
Q_{\text{tỏa}}=\frac{504}{0.70}=\SI{720}{\mega\joule}=\SI{720000}{\kilo\joule}.
\]
Số mol methane:
\[
n_{\ce{CH4}}=\frac{720000}{890}\approx\SI{808.99}{\mole}.
\]
Khối lượng methane:
\[
m=808.99\times16\approx\SI{12944}{\gram}\approx\SI{12.94}{\kilo\gram}.
\]

\textbf{Vậy,} cần khoảng \SI{12.94}{\kilo\gram} methane.
\end{loigiai}

\subsection{Năng lượng trong bài toán nhiệt phân}

Một quá trình thu nhiệt có thể được cung cấp năng lượng bằng cách đốt nhiên liệu. Khi đó, cần tính tổng năng lượng mà quá trình thu nhiệt cần rồi tính ngược lượng nhiên liệu.

\begin{vidu}
Để phân hủy hoàn toàn \(1\) mol \ce{CaCO3} cần cung cấp \SI{178}{\kilo\joule}. Một lò nung xử lí \SI{500}{\kilo\gram} \ce{CaCO3} tinh khiết. Nhiệt được cung cấp bằng một loại nhiên liệu có nhiệt trị \SI{40}{\mega\joule\per\kilo\gram}; hiệu suất truyền nhiệt vào lò là \(50\%\). Tính khối lượng nhiên liệu tối thiểu cần dùng. Cho \(M_{\ce{CaCO3}}=100\) g/mol.
\end{vidu}
\begin{loigiai}
Số mol \ce{CaCO3}:
\[
n=\frac{500000}{100}=\SI{5000}{\mole}.
\]
Năng lượng mà phản ứng phân hủy cần:
\[
Q_{\text{ích}}=5000\times178=\SI{890000}{\kilo\joule}=\SI{890}{\mega\joule}.
\]
Do hiệu suất truyền nhiệt là \(50\%\):
\[
Q_{\text{nhiên liệu}}=\frac{890}{0.50}=\SI{1780}{\mega\joule}.
\]
Khối lượng nhiên liệu:
\[
m=\frac{1780}{40}=\SI{44.5}{\kilo\gram}.
\]

\textbf{Vậy,} cần tối thiểu \SI{44.5}{\kilo\gram} nhiên liệu.
\end{loigiai}

\subsection{Bài toán hỗn hợp nhiên liệu}

Nếu nhiên liệu gồm nhiều chất, năng lượng toàn phần bằng tổng năng lượng do từng thành phần cung cấp:
\[
Q_{\text{tổng}}=Q_1+Q_2+\cdots.
\]
Cần đặc biệt chú ý xem thành phần hỗn hợp được cho theo \textbf{mol, thể tích hay khối lượng}.

\begin{vidu}
Một hỗn hợp khí gồm \(40\%\) propane và \(60\%\) butane theo số mol. Khi đốt hoàn toàn, \(1\) mol propane và \(1\) mol butane lần lượt tỏa \SI{2220}{\kilo\joule} và \SI{2850}{\kilo\joule}. Tính nhiệt lượng tỏa ra khi đốt cháy \(100\) mol hỗn hợp.
\source{Phỏng theo dạng bài LPG trong đề tuyển sinh chuyên Hóa Điện Biên 2026 và Lâm Đồng 2026.}
\end{vidu}
\begin{loigiai}
Trong \(100\) mol hỗn hợp có \(40\) mol propane và \(60\) mol butane. Do đó:
\[
Q=40\times2220+60\times2850=\SI{259800}{\kilo\joule}
=\SI{259.8}{\mega\joule}.
\]

\textbf{Vậy,} \(100\) mol hỗn hợp LPG tỏa ra \SI{259.8}{\mega\joule}.
\end{loigiai}

\begin{tuluyen}
Một thiết bị cần \SI{96}{\mega\joule} năng lượng có ích. Thiết bị sử dụng nhiên liệu có nhiệt trị \SI{32}{\mega\joule\per\kilo\gram} và hiệu suất \(75\%\). Tính khối lượng nhiên liệu cần dùng.
\dapso{\SI{4.0}{\kilo\gram}.}
\begin{loigiaiBT}
\[
Q_{\text{tỏa}}=\frac{96}{0.75}=\SI{128}{\mega\joule},
\qquad
m=\frac{128}{32}=\SI{4.0}{\kilo\gram}.
\]
\textbf{Vậy,} cần \SI{4.0}{\kilo\gram} nhiên liệu.
\end{loigiaiBT}
\end{tuluyen}

\begin{tuluyen}
Đốt cháy hoàn toàn \SI{11.5}{\kilo\gram} ethanol, biết \(1\) mol ethanol tỏa \SI{1370}{\kilo\joule}. Nếu chỉ \(35\%\) năng lượng được chuyển thành công có ích, tính năng lượng có ích thu được. Cho \(M_{\ce{C2H5OH}}=46\) g/mol.
\dapso{\(\approx\SI{119.9}{\mega\joule}\).}
\begin{loigiaiBT}
\[
n=\frac{11500}{46}=\SI{250}{\mole}.
\]
Nhiệt lượng tỏa ra:
\[
Q=250\times1370=\SI{342.5}{\mega\joule}.
\]
Năng lượng có ích:
\[
Q_{\text{ích}}=342.5\times0.35\approx\SI{119.9}{\mega\joule}.
\]
\textbf{Vậy,} năng lượng có ích thu được khoảng \SI{119.9}{\mega\joule}.
\end{loigiaiBT}
\end{tuluyen}
